\reportheading{1. Introduction and design basis}

\subhead{1.1 Design objective}

This project addresses the lighting needs of an open-plan office used for reading, paperwork, and computer work. To support these tasks, a maintained average illuminance of 500 lx is required on each desk, with a uniformity of at least 0.60 and a maximum UGR of 19 [1]. Ultimately, the aim is to provide enough light for these activities without large differences across the working area while reducing excessive glare from the fittings.

\subhead{1.2 Office geometry and proposed use}

The L-shaped office wraps around the central service core. For the BEC lighting power density calculation, the internal floor area of 1042.10 m² is measured to the inner faces of the enclosing walls and includes the column footprints [11].

Within this layout, the office space is arranged with 28 four-person bench groups, totalling 112 desks, each 1.6 by 0.8 m. Their positions serve as reference points for evaluating lighting across the workplace. Accordingly, the calculations below assume a fixed furniture arrangement, meaning a revision would be necessary if the desks are moved significantly.

\subhead{1.3 Room conditions}

Table 1 indicates the key assumptions implemented in DIALux. The relatively light ceiling and walls contribute reflected light onto the desks, whereas daylight is excluded because the window dimensions are not taken into account. Section 7 further evaluates the impact of changing these conditions.

\tabletitle{Table 1. Principal room and calculation assumptions.}

\begin{reporttable}{L{0.5}L{0.5}}
\rowcolor{tablegrey} \textbf{Parameter} & \textbf{Adopted value} \\ \hline
Ceiling / desk / seated eye height & 2.8 / 0.8 / 1.2 m \\
Ceiling / wall / floor reflectance & 80 / 70 / 20\% \\
Column reflectance / maintenance factor & 40\% / 0.80 \\
Furniture and daylight & Proposed desks; no high partitions or simulated daylight \\
\end{reporttable}
